% !TeX root = ../../main.tex
\subsection{Two PCS queries defeat block-local metadata swaps}\label{note:zk-flock-pair-opening}

\paragraph{Conclusion.} Two initial PCS queries defeat the original adjacent metadata swaps, and whole-row shuffling alone retains a second leak. Independent valid counter-label swaps repair it. This note's strongest theorem adds a shared-frame program bank and jointly covers all $37$ BLAKE2s terminal claims, the established Flock endpoints, one selected lane-$59$ query pair and the three computed GKR leaf evaluations, below $2^{-154}$. The memory envelope and shared bus root join below $2^{-151}$. The subsequent replacement in \noteref{zk-flock-children} covers the twelve transmitted children too, without a global shuffle. All these sources use unchanged public heights and verifier equations.

\subsubsection{The encoder reveals sums inside a coefficient block}

Write one initial-lane polynomial as $f(X)=\sum_i f_i\novel_i(X)$. Since $\novel_1(X)=X$ and every higher subspace factor is invariant under $X\mapsto X+1$, two distinct queries $x,x+1$ reveal
\[
 x f(x)+(1+x)f(x+1)
 =\sum_j(f_{2j}+f_{2j+1})\novel_{2j}(x)
 =\zkBlockCollapse{2}(f)(x).
\]
Every exchange of two adjacent coefficients preserves this observable. Equivalently, its difference polynomial is $(1+X)G(X)$ with $G(X+1)=G(X)$, so the displayed combination cancels it exactly. This is a relation among raw authenticated initial leaves; later folding and sumcheck masking cannot erase it.

More generally, for $B=2^k$ define
\[
 \zkBlockCollapse{B}(f)(X)
 =\sum_j\left(\sum_{a=0}^{B-1}f_{Bj+a}\right)\novel_{Bj}(X).
\]
The $B$ queries in any coset $x+\subsp_k$ determine this value. Indeed, write $f=\sum_{a<B}\novel_a F_a$, where each $F_a$ is constant on that coset. The evaluation matrix $(\novel_a(x+s))_{s\in\subsp_k,a<B}$ is invertible by ordinary interpolation, so the verifier recovers all $F_a(x)$ and sums them. The collapse is constant on the coset. If it is nonzero of degree $d$, it vanishes on at most $d/B$ whole cosets. Thus any source preserving all aligned size-$B$ coefficient-block sums has a small-query invariant, regardless of its entropy within those blocks.

\subsubsection{Where the actual stack exposes the invariant}

At the current layout $(\kmem,\mu)=(25,28)$ with table heights $(19,19,19,19,20,18)$ and code height $11$, initial lane $59$ contains five consecutive BLAKE2s count columns, followed by the bytecode-count column and zeros. Its first five size-$2^{18}$ blocks are
\[
 (\mathsf{cnt\_cv1},\mathsf{cnt\_out0},\mathsf{cnt\_out1},
   \mathsf{cnt\_md},\mathsf{cnt\_bc}).
\]
The reference placement algorithm certifies these exact offsets. All swaps in \noteref{zk-flock-columns} exchange physical row indices differing only in bit zero, in every affected column. The compression-input randomness changes none of these counters, and the swaps preserve the bytecode-count table. Hence the entire random contribution to this lane is annihilated by $\zkBlockCollapse{2}$. Unused-memory masks occupy other lanes.

\paragraph{A valid private alias choice.} A fixed public program performs four identical compressions in a frame with integer base $a$, changes frame, performs four more identical compressions, then returns to its public final state. Its private destination has integer base either $a+32$ or $a$, encoded by pointer $\gen^{a+32}$ or $\gen^a$. Both choices are legal: message, chaining value and metadata agree, and repeated writes produce the same digest. Distinct control offsets avoid a write-once conflict when the second frame aliases the first. Each compression has its own code location, so its bytecode read count is one in both executions.

For each of the first four count columns in lane $59$, the first group has counts $1,\gen,\gen^2,\gen^3$. The second group has those same counts with distinct frames, or $\gen^4,\ldots,\gen^7$ with aliasing. Place these eight real rows at logical indices $96\cdot2048+i$, $0\le i<8$. Their physical row indices are $12288+i$. Keep all other rows and padding schedules fixed, with disjoint memory and code, and let $\Delta f$ be the difference of the two lane polynomials. Its nonzero coefficients are exactly
\[
 (\Delta f)_{j2^{18}+12292+i}=\gen^i(1+\gen^4),
 \qquad 0\le j<4,\quad 0\le i<4.
\]
The remaining bytecode-count block does not change. The adjacent collapse is nonzero of degree $798726$; its highest coefficient is $(\gen^2+\gen^3)(1+\gen^4)\ne0$. For $B=4,8,16$, its degree is respectively $798724,798720,798720$, again with nonzero leading coefficient.

A fixed valid completion exists for this obstruction. Fill the other BLAKE2s positions with closed compression/return cycles on disjoint frames, respecting the prescribed random libraries. Even a distinct $32$-cell frame for each compression fits the payload. Fill each of the four other non-JUMP tables through fixed zero-operation blocks of at most sixteen rows and one return, requiring at most $4\cdot2^{15}$ additional JUMPs. Together with the compression returns and real JUMPs this is below $2^{20}$; fixed self-loops supply the remainder. Repeated dummy operations share a few separate frames, and their code blocks fit the public height $11$. The same schedules apply to both real executions. This completion deliberately does not assert any of the separate count-normalization or full-GKR theorems.

The native example \path{crates/lean_vm/examples/zk_flock_pair_opening_audit.rs} verifies two complete proofs of this small-frame program with the same public input and heights. Each execution has counts $(1,1,8,1,4,8)$, so no native fillers complicate the witness pair. It separately checks the two-query identity through the native NTT and authenticated leaves in a small stack. This validates the execution and encoder mechanisms, not a full size-$28$ ZK-mode proof or a Fiat-Shamir probability transfer.

\subsubsection{A quantitative honest-query obstruction}

\begin{theorem}[Block-local metadata privacy floor]\label{thm:zk-flock-pair-floor}
Retain the stated real counted rows at their prescribed positions and use the fixed completion above. Suppose all private padding choices preserve the size-$B$ block sums of lane $59$, where $B\in\{2,4,8,16\}$. Let $D=2^{22+\rexp}$, let $q$ be the actual initial query count, and let $d$ be the nonzero private collapse degree above. Put
\[
 G_B=(D-d)/B,\qquad
 p_j=\sum_{a=0}^j(-1)^a\binom ja\left(\frac{D-a}{D}\right)^q.
\]
Every common simulator for the complete honest-interactive view incurs error at least
\[
 \frac12\left(G_Bp_B-\binom{G_B}{2}p_{2B}\right).
\]
In particular, the adjacent-swap construction cannot reach the statistical target, even with arbitrary compression-input randomness and unused-memory masks.
\end{theorem}
\begin{proof}
There are at least $G_B$ disjoint cosets on which the two fixed collapsed polynomials differ. Whenever every point of one such coset is queried, interpolation recovers distinct values unaffected by the padding choices. Inclusion-exclusion gives $p_B$ for hitting one full coset and $p_{2B}$ for hitting two. Bonferroni gives the public event probability in parentheses. On this event the two supports are disjoint, so their distance is at least that probability, and a common simulator pays at least half. Additional transcript coordinates cannot decrease the distinguishing advantage.
\end{proof}

Exact rational evaluation gives these strict simulator-error lower bounds:
\[
\begin{array}{c|rrrr}
 B\backslash\rexp&1&2&3&4\\\hline
 2&2^{-10}&2^{-13}&2^{-15}&2^{-17}\\
 4&2^{-41}&2^{-48}&2^{-54}&2^{-58}\\
 8&2^{-103}&2^{-118}&2^{-130}&2^{-141}\\
 16&2^{-226}&2^{-258}&2^{-283}&2^{-306}
\end{array}
\]
The last line does not prove that blocks of sixteen are safe; it only ceases to exclude $2^{-128}$ by this particular event. The theorem concerns the specified execution-preserving completion and block-preserving randomizers. It is not a claim against every possible count-normalized completion, arbitrary row permutation, additional metadata source, or a prover allowed to canonicalize the real execution instead of retaining its counted rows. Its purpose is to reject the presently proved terminal-only source as a sufficient full-view construction.

\subsubsection{A compatible source that is not block-local}

Keep the original library and its independent bits. Add a second copy at disjoint complementary positions, with code locations shifted by $64$ and integer frame bases shifted by $2^{22}$ into the next payload interval. All $16256$ compressions in this second copy use the same fixed valid compression. Uniformly permute their complete BLAKE2s rows across the new positions, leaving their closing JUMPs in place. Whole-row permutation preserves every bus multiset, memory image and count-column product. Equal compression values make the packed Flock witness invariant pointwise.

The new placement is explicit. For program rows map old logical banks $(0,1,2,4,8,16,32)$ to $(5,6,7,9,10,11,12)$ without changing $s$. For frame-bank rows map banks $64+(0,1,2,4,8,16)$ to $64+(3,5,6,7,9,10)$. For count-bank rows keep the bank and change the top-three-bit within-bank sector from $3$ to $6$. These positions avoid the original library, the mandatory three-point support and the five general metadata rows. All remain in the first $96$ banks. The new global permutation crosses the old coefficient blocks.

\begin{theorem}[Whole-row shuffling retains a cross-column leak]\label{thm:zk-flock-shared-count-floor}
Use the original adjacent-swap source and the second library with whole-row shuffling only. For an execution-preserving fixed completion, a valid private read-target choice gives common-simulator error strictly greater than $2^{-15},2^{-19},2^{-22},2^{-25}$ at rates one through four.
\end{theorem}
\begin{proof}
Each second-library row has all nine memory-count columns equal: they are all $1$ in a fresh-frame row, or all $1$ or $\gen$ in a repeated-pair row. Whole-row shuffling preserves their pointwise equality. Consequently its change to lane $59$ has the form
\[
 (1+\novel_{2^{18}})(1+\novel_{2^{19}})S
   +\novel_{2^{20}}T,
 \qquad \deg S,\deg T<2^{18}.
\]
On the coset $\gen^{18}+\subsp_{18}$ the three displayed novel factors are respectively $1,0,0$, so this change vanishes identically. The coset is closed under $x\mapsto x+1$; the original library's contribution is still canceled by $\zkBlockCollapse{2}$.

A fixed small-frame program first uses a private pointer to read either its chaining-value cell at offset $9$ or an identical decoy at offset $20$. It then performs eight identical compressions and returns to the public final state. Hint the chaining value and pointer, so initializing the chaining value does not itself consume a read label. In the first execution the earlier DEREF advances only $\mathsf{cnt\_cv1}$ among the selected BLAKE2s columns. With the same real-row placement as above, the lane difference is exactly
\[
 (\Delta f)_{12288+i}=\gen^i(1+\gen),\qquad 0\le i<8.
\]
All other lane-$59$ coefficients agree. Its adjacent collapse has degree $12294$ and leading coefficient $\gen^6(1+\gen)^2\ne0$. Thus at least $(2^{18}-12294)/2=124925$ adjacent pairs in this coset distinguish the witnesses, independently of both libraries' coins. The same Bonferroni bound as Theorem~\ref{thm:zk-flock-pair-floor}, with $G_2=124925$, gives the stated exact rational floors.
\end{proof}

The native example also verifies this second witness pair, with common public input, outputs and counts $(1,1,8,1,2,8)$. Both executions use one $32$-cell frame. For the full-layout obstruction, relocate the real frame into a disjoint payload slot, for example at integer base $3\cdot2^{22}-128$, reserving it from other fillers; the isolated native example uses base $1280$. Hence ordinary small-frame execution does not itself hide read-access patterns. The theorem retains the real counted rows and uses the fixed completion; it does not exclude arbitrary count-normalized completions or other padding sources.

\paragraph{Independent labels repair this source-level defect.} For every one of the second library's $4096$ repeated pairs, independently choose whether to exchange the two labels in each of its ten count columns, then globally shuffle its $16256$ BLAKE2s rows. These are $40960$ independent fair bits in addition to the permutation. The two rows in a pair have identical instruction, addresses and values; each exchange preserves the full address-specific read chain, count-column product and bus multiset. Memory, code and every compression value remain fixed. The nine memory-count columns need no longer agree at each row. A common swap bit across columns would preserve their equality and would not repair the theorem's obstruction.

\begin{proposition}[Preserving the established joint theorem]\label{prop:zk-flock-global-extension}
The added independent library leaves the error bounds in Theorem~\ref{thm:zk-flock-all-columns} and Corollary~\ref{cor:zk-flock-memory-interface} unchanged, under their common-completion assumptions.
\end{proposition}
\begin{proof}
Condition on the entire added permutation and its independent counter-label bits. Its compression values and all other rows are a fixed valid complement for the original library's terminal proof. Disjoint code and memory preserve that library's directions and read chains. Its independent bits therefore give the same witness-free target law with the same uniform error bound. The Flock witness does not depend on any added choice, so the existing endpoint proof applies unchanged. Average over those choices and use the same unused-memory root/envelope composition.
\end{proof}

This uses $32512$ of the $98304$ complementary compression positions, leaving $65792$ positions for independent general inputs. Both count banks share a frame within each repeated pair, so Flock needs at most $188416$ frames of $32$ cells, or $6029312$ cells, leaving $25585$ payload slots of size $256$. The two libraries use $72$ code locations. No table height or verifier equation changes. Common allocation with other libraries remains open.

\subsubsection{A quantified repair for the exposed query pair}

\begin{lemma}[Random matching fills a binary image]\label{lem:zk-random-matching}
Let $N$ vertex values lie in $\K$, with $t$ disjoint affine bases over $\Ftwo$. Choose a uniform matching of $M$ disjoint edges, leaving $U=N-2M$ vertices unmatched. With independent fair edge bits, the sum of the selected endpoint differences is within distance
\[
 \zkMatchingError
 =(2^{64}-1)\max_{t\le m\le N-t}
 \frac{[z^m](1+z^2)^M(1+z)^U}{\binom Nm}
\]
of uniform $\K$, even jointly with the matching.
\end{lemma}
\begin{proof}
Conditional on the matching, the sum is uniform on the binary span of its endpoint differences. A nonzero binary linear functional annihilates this span exactly when all edges are monochromatic under its two-coloring of the vertices. Every affine basis contains both colors, so the number $m$ of one-colored vertices lies between $t$ and $N-t$. For a uniform matching, place these $m$ colors uniformly among its fixed edge and singleton slots. The displayed coefficient counts the assignments with zero or two colored endpoints on every edge, out of $\binom Nm$ total assignments. Union-bound over the $2^{64}-1$ nonzero functionals. Outside that event the conditional distribution is exactly uniform, proving the joint claim.
\end{proof}

\begin{theorem}[Hiding the shared-count attack projection]\label{thm:zk-flock-matching-projection}
For the repaired library and any fixed $x\in\gen^{18}+\subsp_{18}$, the observable $\zkBlockCollapse{2}(f)(x)$ from lane $59$ is within $2^{-160}$ of uniform $\K$, for every fixed valid completion outside the added randomizers. This includes joint disclosure of the added permutation, but not other observations of its counter-label bits. The claim uses the exhaustive native-field certificate described below.
\end{theorem}
\begin{proof}
Let $P$ be the $16256$ public positions of the second library and give position $s$ the value $\novel_{s-(s\bmod2)}(x)$. Its $4096$ repeated count pairs become a uniform matching on $P$ after the global shuffle. On the specified coset, the first four lane-$59$ selector weights are $(1,1,0,0)$ and the bytecode selector is zero. For each transported pair, the independent swap bits for $\mathsf{cnt\_cv1}$ and $\mathsf{cnt\_out0}$ therefore contribute their fair XOR times $(1+\gen)$ times the endpoint difference. The common whole-row part vanishes, and all remaining contributions are a fixed translation. Multiplication by $1+\gen\ne0$ preserves binary rank.

For every one of the $131072$ adjacent pairs in this coset, the certificate finds $112$ disjoint affine bases among the $8128$ even positions. Each even position has its odd neighbor in $P$ with the same collapsed value, so their copies give $t=224$ disjoint affine bases among all vertices. Apply Lemma~\ref{lem:zk-random-matching} with $(N,M,U)=(16256,4096,8064)$. Exact integer arithmetic gives
\[
 \max_{224\le m\le N-224}
 \frac{[z^m](1+z^2)^{4096}(1+z)^{8064}}{\binom Nm}
 =\frac{[z^{224}](1+z^2)^{4096}(1+z)^{8064}}{\binom N{224}},
 \qquad \zkMatchingError<2^{-160}.
\]
This proves the stated bound uniformly over $x$ and the fixed translated witness contribution.
\end{proof}

\begin{corollary}[The two raw lane values]\label{cor:zk-flock-matching-pair}
For any fixed pair $x,x+1$ in the same coset, the joint lane-$59$ values $(f(x),f(x+1))$ are within $2^{-160}$ of uniform $\K^2$. This uses both the added source and the original library's independent counter-label bits. All other padding choices may be fixed. The claim can retain the added permutation, but not other observations of the masking bits.
\end{corollary}
\begin{proof}
Use the original $448$ independent $\mathsf{cnt\_cv1}$ swap bits, holding every other original bit fixed. For an even pair position $s$, its contribution to the two evaluations is
\[
 (1+\gen)\novel_s(x)\,(1+x,x).
\]
The exhaustive certificate verifies that these $448$ scalars span $\K$ over $\Ftwo$, for every $x$ in the coset. Thus the original bits supply a perfectly uniform scalar in the kernel of $(u,v)\mapsto xu+(1+x)v$, independently of the entire added source. Theorem~\ref{thm:zk-flock-matching-projection} hides the quotient coordinate below $2^{-160}$. Since $x\ne0$, the map $(u,v)\mapsto(xu+(1+x)v,v)$ is invertible. Conditional uniformity of its second coordinate and approximate uniformity of its first prove the joint claim.
\end{proof}

\paragraph{Reproducible exhaustive certificate.} The audit's \texttt{--matching} option passes the exact sorted even source positions, shuffled once with fixed seed $467$ solely for the basis-search order, to the native example. It enumerates every even $x$ in the coset and greedily extracts $112$ disjoint affine bases using binary Gaussian elimination on actual $\K$ values. It also certifies rank $64$ for the original $448$ adjacent-swap scalars at every query. Odd $x$ give the same collapsed values. At every query it cross-checks a sparse polynomial against the native additive NTT. This is exhaustive finite-field evidence, not sampled rank extrapolation or a Lean formalization. The arithmetic bound uses the coefficients $a_m=[z^m](1+z^2)^M(1+z)^U$, computed exactly by
\[
 (m+1)a_{m+1}
 =(U-m)a_m+(2M-m+1)a_{m-1}+(N-m+2)a_{m-2},
 \qquad a_0=1,
\]
with negative indices zero. Differentiating the generating polynomial proves the recurrence. The audit checks symmetry, cross-checks selected coefficients with the binomial sum, and maximizes the exact ratios over the entire integer interval; it assumes no monotonicity.

\paragraph{Scope of the isolated repair.} The corollary also simulates the following marginal: draw all honest initial query indices, select the first complete pair in this coset by a public fixed rule, and disclose its two lane-$59$ values, or a fixed marker if there is no pair. The simulator uses two independent uniform field values, with error below $2^{-160}$; no union bound over possible selected pairs is needed. Separately, a union bound gives a single bad-permutation event below $2^{-143}$ outside which every pair's scalar map has full rank. This does not make those many scalars jointly independent. Preserving the earlier terminal theorem as a marginal does not by itself prove its joint composition with this opening pair; the next theorem accounts for the shared bits.

\subsubsection{Joint composition with the terminal and memory interfaces}

\begin{theorem}[Terminal interface with one actual query pair]\label{thm:zk-flock-pair-interface}
Use the two libraries above and the common-completion assumptions of Theorem~\ref{thm:zk-flock-all-columns}. Include all honest coins and, when present, the first complete query pair in $\gen^{18}+\subsp_{18}$ under a fixed public ordering. Its two lane-$59$ values, all $37$ BLAKE2s terminal evaluations and the established Flock endpoints have a joint witness-free simulator with error at most
\[
 \zkTripleError+3\delta_{\mathrm{span}}+\zkFixedError{64}
 +\frac{4283}{2^{192}}+\zkMatchingError<2^{-154}.
\]
The same statement holds for any one fixed pair in this coset. This is an enlarged algebraic projection: all other query values, omitted table/Flock messages, GKR and authentication remain excluded.
\end{theorem}
\begin{proof}
Write the two lane values as $(u,v)$ and use coordinates $(xu+(1+x)v,v)$. Every original adjacent swap preserves the first coordinate. Condition on the complete added library's permutation and label bits, all compression values and all other private coins except the original metadata bits. Its first coordinate is then fixed.

The original $448$ $\mathsf{cnt\_cv1}$ bits contribute to its terminal evaluation a common nonzero selector times the twelve-coordinate equality weights. They contribute to $v$ the fixed $64$-bit linear map with column weights
\[
 (1+\gen)x\novel_{\zkFlockOrder(\ell_i)}(x),\qquad i\in S,
\]
where $\ell_i$ are the count-bank pair positions. This map is onto $\K$ for every $x$ in the coset, by the exhaustive kernel-rank certificate. It depends on $x$ but not on the honest table point. Theorem~\ref{thm:fixed-observation-span} therefore makes the terminal $\mathsf{cnt\_cv1}$ evaluation and $v$ jointly uniform on good table coins. No new source bits are required.

Use these bits last in the earlier triangular argument. The eight program directions, frame direction, and other nine counter directions make their respective terminal coordinates uniform first; their changes to $v$ are absorbed by the final joint replacement. The mixed span also supplies the ordinary count span for the other eight memory-count columns. The remaining span events are the ordinary program grid and the two geometrically weighted frame/bytecode grids. Including their selectors, the complete metadata replacement costs
\[
 3\delta_{\mathrm{span}}+\frac{96}{|\E|}+\zkFixedError{64}.
\]
It yields nineteen independent uniform extension elements and an independent uniform $v$, while retaining the actual collapse, Flock values and endpoints, and the added library's full randomness.

Now replace the collapse using Theorem~\ref{thm:zk-flock-matching-projection}, at cost $\zkMatchingError$. The first replacement was uniform in every fixed added choice, so this second hybrid does not condition a marginal bound on disclosed terminal values. The added randomizers preserve the Flock witness pointwise. Apply the value/endpoint simulator of Theorem~\ref{thm:zk-flock-metadata-endpoints}, costing $\zkTripleError+4187/|\E|$, and invert the two-coordinate transformation. All bounds are uniform in $x$; its public selection uses only independent honest query coins. If there is no pair, use the earlier terminal simulator and a fixed marker. Adding the three costs proves the formula.
\end{proof}

\begin{corollary}[Memory, bus root and the selected opening pair]\label{cor:zk-flock-pair-memory}
Under the same layout, reservation and common-completion assumptions as Corollary~\ref{cor:zk-flock-memory-interface}, its joint view can also retain the selected lane-$59$ query pair, with total error below $2^{-151}$.
\end{corollary}
\begin{proof}
Repeat that corollary's hybrids: replace the shared root inside the unused-memory kernel, then replace the three-limb memory envelope conditional on all instruction data. Lane $59$ contains no memory-value coefficients, so these replacements retain the two new values exactly. Finally apply Theorem~\ref{thm:zk-flock-pair-interface}, not the former marginal terminal theorem. The exact bound is
\[
 \varepsilon_{\mathrm{decouple}}+\frac{3(22+12)}{2^{192}}
 +\zkTripleError+3\delta_{\mathrm{span}}+\zkFixedError{64}
 +\frac{4283}{2^{192}}+\zkMatchingError<2^{-151}.
\]
\end{proof}

\subsubsection{Retaining the actual GKR count interface}

Use the denser support $\zkDenseSupport$ of Theorem~\ref{thm:dense-fixed-observation-span} for each of the original library's ten counter-label directions. These $1280$ pairs already exist inside its $4096$-pair count bank. This adds $10(1280-448)=8320$ independent fair bits and no rows, frames, code locations or committed coefficients. Every choice still preserves read chains, count products, memory and the packed Flock witness. The previous source is a subset; conditioning on the added bits preserves its existing theorems.

Let $\zeta$ be the honest GKR leaf point. At this layout the native push/pull cube has depth $26$, with the depth-$24$ count prefix padded to the same cube. Define the auxiliary vector
\[
 \zkLeafCounts=
 \bigl(\mle T_c(\zeta_{<18})\bigr)_{
 c\in\{\mathsf{cnt\_m0},\ldots,\mathsf{cnt\_m3},
 \mathsf{cnt\_cv0},\mathsf{cnt\_cv1},
 \mathsf{cnt\_out0},\mathsf{cnt\_out1},\mathsf{cnt\_md},\mathsf{cnt\_bc}\}}.
\]
Here $T_c$ is the physical BLAKE2s count column. These ten evaluations are not sent individually: \texttt{tables\_and\_prods\_at} computes them inside the native decomposition, and their weighted combinations enter its table forms. Retaining them is an explicit enlargement of the proof interface.

\begin{theorem}[Count interface, terminal claims and an opening pair]\label{thm:zk-flock-count-interface}
Use the denser counter source and the same common-completion assumptions. The view of Theorem~\ref{thm:zk-flock-pair-interface} can also retain $\zkLeafCounts$, with joint error at most
\[
 \begin{aligned}
 &\zkTripleError+3\delta_{\mathrm{span}}+\zkFixedError{64}
 +\zkDenseError{256}+\zkDenseError{192}\\
 &\hspace{2em}+\frac{4317}{2^{192}}+\zkMatchingError
 <2^{-154}.
 \end{aligned}
\]
With the three-limb memory envelope and shared bus root, the combined exact bound is still below $2^{-151}$. Earlier GKR messages and the count root remain excluded.
\end{theorem}
\begin{proof}
The GKR row point and the subsequent table terminal point are independent honest extension points; apply the physical-to-logical coordinate permutation to each. Fix a selected query $x$, and use the original $448$-position subset first only to certify a retained map's rank. Theorem~\ref{thm:fixed-observation-span} makes the $\mathsf{cnt\_cv1}$ evaluation at $\zeta$ and its raw query-kernel coordinate a full $256$-bit map, except with probability $\zkFixedError{64}$. The exhaustive query certificate supplies the fixed $64$-bit rank for every $x$ in the coset. Consequently every other memory-count evaluation at $\zeta$ also has full rank on this subset. The weighted bytecode-count evaluation at $\zeta$ has full rank except with probability $\delta_{\mathrm{span}}+36/|\E|$. Passing to the denser source cannot reduce these ranks.

Condition on $\zeta$ and $x$. Apply the dense fixed-disclosure theorem at the independent table terminal point: retain $256$ bits for $\mathsf{cnt\_cv1}$, and $192$ bits for each other count. The former event already contains the common unweighted two-point span for the other eight memory-count columns. The bytecode column needs its own dense event and geometric transformation, costing $\zkDenseError{192}+24/|\E|$. Thus each memory-count bit family jointly supplies its two extension evaluations, with the additional independent query coordinate from $\mathsf{cnt\_cv1}$; the bytecode family supplies its two evaluations too.

Insert these count replacements last in the program/frame triangular proof. The two unchanged program/frame spans cost $2\delta_{\mathrm{span}}+48/|\E|$. Selectors cost at most $12/|\E|$ at the table point and $10/|\E|$ at the GKR point. Together with the earlier bytecode span and its geometric transformation, these field terms sum to $(48+36+24+12+10)/|\E|=130/|\E|$. The replacement yields nineteen terminal metadata uniforms, ten independent count-interface uniforms and an independent raw query-kernel uniform, retaining all compression data and the added library's full randomness.

Now perform the same quotient replacement and value/endpoint simulation as before, adding $\zkMatchingError+\zkTripleError+4187/|\E|$. If no pair is queried, prove the enlarged statement using the fixed auxiliary pair at $\gen^{18}$ and then discard its two values. This gives the same bound without a separate fallback case. Finally, the memory/root hybrids retain $\zkLeafCounts$ exactly because unused-memory masks change no table count. Adding their existing exact bound proves the second assertion.
\end{proof}

\paragraph{What this controls in the actual bus.} Let the three native GKR leaf evaluations be $(P,Q,C)$. Extract from the actual table forms their coefficients on the ten BLAKE2s count columns. This gives a public $3\times10$ matrix mapping $\zkLeafCounts$ to the entire count-dependent part of $(P,Q,C)$. No count appears in a quadratic monomial. Its first row is $\gen$ times its second: a counted push increments a label and its corresponding pull does not, with identical block placement. Thus counter swaps leave $P+\gen Q$ fixed.

The matrix generically has rank two. For $\mathsf{cnt\_cv1}$ and $\mathsf{cnt\_out0}$, the pull block indices above the $18$ row bits are $191,192$, while the count block indices are $49,50$. Their $2\times2$ minor is the common count-coordinate fingerprint weight times
\[
 \eq(\zeta_{18:26},191)\eq(\zeta_{18:26},50)
 +\eq(\zeta_{18:26},192)\eq(\zeta_{18:26},49).
\]
The two products have different factorizations into coordinate factors and their complements, so this degree-at-most-$16$ polynomial is nonzero. The minor vanishes with probability at most $20/|\E|$, including the fingerprint weight. No extra error is needed to simulate the count contribution: the simulator can apply the actual matrix to the simulated vector even on degenerate coins. But the residual address/value contribution to $P+\gen Q$ remains an essential independent obligation; a count-only extension cannot hide it.

\subsubsection{A same-frame source for the full leaf triple}

Counter swaps cannot move the residual $\zkLeafResidual=P+\gen Q$. A small change to the program bank supplies this missing direction. This is a new source variant; the preceding theorems retain their original source and allocation.

\paragraph{Valid source and placement.} In program bank zero, replace $S$ by $\zkSharedPcSupport=\zkDenseSupport\setminus\{192,193\}$, of size $1278$. Each pair now uses one common $32$-cell frame, the same compression and baseline operands, at code locations $1024,1026$. Its closing JUMPs use control offsets $(16,17,18)$ and $(24,25,26)$ respectively. The two distinct return addresses therefore occupy disjoint cells, and both cycles are valid in write-once memory. The nine compression addresses are each read twice, so their labels are $1,\gen$; the two bytecode labels agree because each code location is used once per pair. Independently exchange the complete BLAKE2s rows at their designated adjacent positions.

The two omitted pair indices are exactly the new collisions with the five existing general-input metadata rows. Keep those rows unchanged. The denser first bank also meets the second library's old program-bank-zero placement; remap that second bank's logical banks $(0,1,2,4,8,16,32)$ to $(13,14,20,21,22,23,24)$. All other banks keep their previous placements. The new positions are still adjacent pairs in the complement of the mandatory binary source, within the first $96$ banks. The exhaustive query certificate is rerun on this changed second library, rather than presumed to survive relocation.

\paragraph{The actual residual direction.} Let $A_P,A_Q,A_C$ be the three BLAKE2s table forms obtained from the native depth-$26$ bus layout, including their block selectors and fingerprint weights. The residual form is $A_P+\gen A_Q$. Its state-PC term cancels, every counted-label term cancels, and its only remaining PC term is
\[
 (1+\gen)\eq(\alpha,1)\eq(\zeta_{18:26},185)\,\mathsf{pc},
\]
where $185$ is the bytecode block index above the row bits. In a same-frame pair every memory address, operand and compression value agrees. Consequently the residual change is exactly this coefficient times $\gen^{1024}+\gen^{1026}$ times the pair's equality-weight sum. The coefficient is a nonzero polynomial of degree at most $12$ in honest coins. This is a direction in the actual full leaf evaluation: framework terms and other tables are a fixed translation during these swaps.

\begin{lemma}[Two evaluations after two omitted source bits]\label{lem:zk-shared-pc-span}
The shared-PC bank jointly masks its table-PC evaluation and its residual evaluation, except with probability at most
\[
 \delta_{\mathrm{span}}+\zkDenseError{194}+29/|\E|,
\]
in addition to the already charged table selectors.
\end{lemma}
\begin{proof}
The original $448$-position support avoids the omitted indices. At the earlier GKR row point it spans $\E$ except with probability $\delta_{\mathrm{span}}+12/|\E|$. On that event, the map from the full dense source to this evaluation and the two individual bits at indices $192,193$ is onto $\E\times\Ftwo^2$. It is fixed independently of the subsequent table point. Theorem~\ref{thm:dense-fixed-observation-span}, retaining $194$ bits, therefore makes its joint map with the table evaluation onto, at additional cost $\zkDenseError{194}$. Restrict the two distinguished bits to zero. Every fiber still has the full two-evaluation image, so removing these positions incurs no conditioning penalty. The GKR pair selector has five factors and the residual coefficient has degree at most twelve; their failures add at most $17/|\E|$. This proves the bound.
\end{proof}

\begin{theorem}[Full leaf boundary, terminal interface and an opening pair]\label{thm:zk-flock-leaf-interface}
Use this shared-PC variant, the dense counter source and the remapped second library, under the same common-completion assumptions. All $37$ BLAKE2s terminal claims, the established Flock endpoints, one selected lane-$59$ query pair and the three actual GKR leaf evaluations $(P,Q,C)$ have a joint witness-free simulator with error at most
\[
 \begin{aligned}
 &\zkTripleError+4\delta_{\mathrm{span}}+\zkFixedError{64}
 +\zkDenseError{256}+\zkDenseError{192}+\zkDenseError{194}\\
 &\hspace{2em}+\frac{4366}{2^{192}}+\zkMatchingError<2^{-154}.
 \end{aligned}
\]
The full three-limb memory envelope and shared bus root can also be included, with total error below $2^{-151}$. The ten auxiliary count evaluations are not retained in this theorem.
\end{theorem}
\begin{proof}
Fix all compression inputs, unused memory and private choices outside the original metadata source. The seven operand banks first mask their seven terminal operands. The frame bank masks the terminal frame pointer; their other effects are left for subsequent banks. The shared-PC bank then masks the terminal PC and $\zkLeafResidual$ jointly by the lemma, absorbing every earlier contribution to those two coordinates. It changes neither the operand nor frame terminal evaluations.

Next use exactly the dense two-point counter maps established in Theorem~\ref{thm:zk-flock-count-interface}. Conditional on all preceding bank bits, they jointly mask the ten terminal counts and supply ten independent earlier count scalars, plus the raw query-kernel coordinate. Their map into $(P,Q,C)$ has rank two outside the degree-$20$ minor exception. Its image is exactly $\{(p,q,c):p+\gen q=0\}$. Applying it to the uniform count scalars fills the affine plane with the already masked residual. Discard the auxiliary scalars: the resulting full leaf triple is uniform on $\E^3$, independent of the terminal metadata and raw query-kernel coordinate. No assertion is made about retaining both the scalars and their translated leaf image.

Thus the metadata replacement has the earlier count-interface cost plus $\delta_{\mathrm{span}}+\zkDenseError{194}+49/|\E|$: $29/|\E|$ from the new lemma and $20/|\E|$ from the rank-two minor. Every original pair is still physically adjacent and preserves the collapsed query coordinate. In the shared-PC bank the nine memory-count differences agree, so its two raw lane values on this coset are unchanged as well. The remapped second library has the same exhaustive $224$-affine-basis certificate, so its independent matching bits hide the quotient at the same cost $\zkMatchingError$. The value/endpoint replacement is unchanged. This gives the displayed total, with $4317+49=4366$.

For the memory/root extension, perform the metadata and query replacements first. All their row permutations and same-tuple label exchanges preserve the individual push and pull multisets, hence the actual shared product root, not merely bus validity. They also preserve the entire memory image and envelope. The resulting uniform leaf absorbs its original framework translation, even though that translation depended on memory. Now apply the unused-memory root and envelope hybrids, retaining these independent uniforms and the actual compression-value/Flock interface. Finally replace that interface by its earlier simulator. This order avoids assuming that the actual leaf would remain fixed under an unused-memory change. Adding the same exact memory/root error proves the second claim.
\end{proof}

\paragraph{Resources and evidence.} The first library now uses $17916$ BLAKE2s rows and the same number of returns; the second still uses $16256$. Together they occupy $34172$ complementary positions, leaving $64132$ general-input positions, and still use $72$ code locations. Sharing the new $1278$ frames lowers the total Flock allocation to at most $187138$ frames, or $5988416$ cells. The remaining payload has $25744$ slots of size $256$. All public heights and verifier equations are unchanged. \auditref{zk_flock_leaf_audit.py} checks the actual forms, valid cycles, disjoint reservations, native binary ranks $5824$ for the intermediate counter/residual map and $4288$ for the actual leaf/terminal/kernel map, and the exact error ledger. Its \texttt{--matching} option exhaustively certifies the relocated source. The symbolic span proof, not the sampled ranks, supplies the uniform theorem.

\paragraph{Remaining opening and GKR obligations.} The three leaf evaluations are the values computed after interpolating the final GKR children at two fresh coins. This theorem does not include the twelve transmitted children themselves or the preceding $24$ sumcheck rounds of the final depth-$26$ reduction. Nor does it include earlier layers, the count root, table coefficients, middle Flock rounds, other queries or authentication. Extending the joint boundary to those actual messages, and proving one common valid sampler, remain essential. The fixed-disclosure bounds are finite budgets, not permission to condition on an arbitrary transcript.

\paragraph{Evidence and literature boundary.} \auditref{zk_flock_pair_opening_audit.py} checks exact stack positions, coset interpolation, all displayed probabilities, disjoint library allocation and full counted-bus validity before and after the added randomizers. It certifies native binary rank $5632$ for the nineteen terminal columns, ten earlier count evaluations and query-kernel coordinate together. It also extracts the actual depth-$26$ bus forms and checks their counter dependence against local label exchanges. The exact composition ledger, native valid witness pairs and exhaustive query certificate are retained. Flock's unused witness slots are forced to zero, as specified and implemented in \path{crates/flock/src/hash.rs}, so they cannot serve as free masks. Bootle, Chiesa and Liu~\cite{bcl}, Section~2.4, add a separately designed R1CS gadget with controlled randomness and also rerandomize their sumcheck factors; this is not a theorem for arbitrary random compressions in leanVM's unchanged circuit and PCS.
